\section{Inscriptions and the fold}\label{sec:inscriptions}

A frog touches Bitcoin in two kinds of inscription: the mint document and the
crossing state body. Both use ord's envelope grammar. The mint is a
self-contained HTML document; the crossing is 55 bytes of binary state. The
mint envelope in the OP\_RETURN output names the seed and name; the mint
inscription carries the art. The indexer validates both, because the
inscription is the part a stock ord viewer can read without HOPPER software.

\subsection{The mint envelope}

The mint transaction carries one OP\_RETURN output whose push is the
following byte string.

\begin{table}[H]
\centering
\small
\caption{The \code{0x08} mint envelope. The total length is at most 80
bytes. The \code{0x07} version is accepted only for the grandfathered sat.}
\label{tab:mint-envelope}
\begin{tabular}{@{}rrll@{}}
\toprule
\textbf{Offset} & \textbf{Length} & \textbf{Field} & \textbf{Value} \\
\midrule
0  & 3 & magic & \code{"FRG"} \\
3  & 1 & version & \code{0x08} (legacy \code{0x07} only for the grandfather) \\
4  & 1 & action & \code{0x01} (MINT) \\
5  & 4 & seed & \code{u32} little-endian \\
9  & 1 & name length & $0$ to $38$ \\
10 & $n$ & name & UTF-8, $n$ bytes \\
\bottomrule
\end{tabular}
\end{table}

Exactly one \code{FRG} envelope is allowed per transaction. A malformed or
duplicated envelope invalidates the protocol message: the transaction still
moves sats, but it is not a mint. The same seed and name appear in the mint
inscription reveal; the indexer validates both. This redundancy is
intentional. The OP\_RETURN makes the mint cheap to index; the inscription
makes it readable by ord alone.

\subsection{The ord envelope}

The mint reveal's tapleaf contains ord's standard envelope:
\begin{lstlisting}[language={},caption={Mint inscription envelope (grammar).},label={lst:mint-env}]
OP_FALSE OP_IF
  "ord"                              ; 3-byte protocol id
  0x01 "text/html;charset=utf-8"     ; tag 1: content type
  0x00                               ; body tag
  <body, chunked at 520 bytes>       ; the fold document
OP_ENDIF
\end{lstlisting}
The crossing envelope is the same grammar with a binary content type, a
parent tag and a fixed body:
\begin{lstlisting}[language={},caption={Crossing inscription envelope (grammar).},label={lst:cross-env}]
OP_FALSE OP_IF
  "ord"
  0x01 "application/octet-stream"
  0x03 <mint_txid:32 || minimal-LE mint_index>   ; tag 3: parent
  0x00 <55-byte body>
OP_ENDIF
\end{lstlisting}

The parent value is the mint transaction id in internal byte order, followed
by the mint output index. Index zero is encoded as the bare 32-byte txid; a
non-zero index is appended as a minimal little-endian integer. Ord's rule is
strict: a one-to-three-byte index whose last byte is \code{0x00} is
non-minimal and rejected. The parent tag is mandatory on a crossing; it is
what makes \code{/r/children} work and what lets a viewer discover the
crossing history of a mint.

\subsection{The mint body: a self-contained proof-checking fold document}

The mint body is the fold document returned by
\code{live_v8_html} for a \code{(seed, sat, name, gen_hint)} tuple, with
content type \code{text/html;charset=utf-8}. It has five properties:

\begin{enumerate}[leftmargin=1.6em,itemsep=0.25em]
  \item \textbf{No external origin.} There is no \code{http://},
        \code{https://}, bare host, or protocol-relative URL anywhere in the
        document. This is grep-enforced in the implementation.
  \item \textbf{Two inline sprites.} The document embeds two
        \code{data:image/png;base64,...} 16$\times$16 images, one idle and
        one dead.
  \item \textbf{A static name band.} The card carries the Hopper's name, its
        generation badge (\code{G1}, \code{G2}, \ldots) and the road mark with
        an unknown crossing count, for example \code{NAME}
        \textperiodcentered{} \roadglyph{} \code{—}. It never prints a number.
  \item \textbf{An embedded parameter block.} The document contains
        \begin{center}{\scriptsize\code{V8={sat, genHint, name,
        [grandfatheredSat,] p:{g,b,gi,cap,hold,gl,fin}}}}\end{center}
        where \code{g} is the genesis height, \code{b} the base, \code{gi}
        the growth interval, \code{cap} the ceiling, \code{gl} the generation
        length and \code{fin} the finality. \code{grandfatheredSat} is
        emitted only when the compiled value is non-zero.
  \item \textbf{The canonical fold source.} The embedded \code{live_v8.js}
        reads only relative \code{/r/...} and \code{/content/...} endpoints
        on the viewer's own ord node.
\end{enumerate}

There is no \code{STATE FROM THE CHAIN} caption in v8. The static band
states only what the mint transaction itself can prove: the name, the
generation hint from the intended mint height, and that the count is unknown.
The live band is authoritative only when the fold has verified every
body-valid child against the raw transaction that published it. If a proof is
missing or false, the band is \code{NAME} \textperiodcentered{}
\roadglyph{} \code{—} and no count or \code{HOME} is claimed. When the proof
holds:
\begin{center}
\small
\code{alive}:\quad \nameband{ 12}\qquad
\code{dead}:\quad \texttt{RIP }\nameband{ 12}\\[1mm]
\code{home}:\quad \nameband{ 25 HOME}
\end{center}
The count is \code{dist}, the number of completed crossings, and it is shown
for the dead as well. \code{RIP } (with its trailing space) is prepended to
the name only when the fold has settled the frog as \code{dead}. The road
glyph is illustrative; the canonical text form is \nameband{ N}, and that
is what tests compare. The implementation uses the road emoji U+1F6E3
between the name and the count; this paper draws the mark instead, because
the body font carries no emoji glyph.

\code{gen_hint} is the generation of the height the wallet expects the mint
to confirm at. It is a display hint, never consensus: the indexer derives
\code{gen} from the confirmed \code{mint_height} and ignores \code{gen_hint}.
The live fold shows the authoritative \code{G} when it can read the mint
inscription's height and falls back to the hint when it cannot.

The PNG encoder is canonical and deterministic. It emits a PNG signature, an
IHDR chunk (8-bit RGBA, no interlace), exactly one IDAT chunk, and IEND. The
zlib stream begins with \code{78 01}, uses \emph{stored} deflate blocks of at
most 65{,}535 bytes with BFINAL set on the last, and ends with an Adler-32
checksum. Every scanline begins with filter byte zero. CRC-32 is computed by
hand per chunk. Two encoders that follow these rules produce the same bytes
for the same sprite.

\begin{sloppypar}
\paragraph{One accepted digest.}
The indexer accepts exactly one v8 mint-body digest per
$(\mathrm{seed},\mathrm{sat},\mathrm{name},\mathit{gen\_hint})$: the SHA-256
of \code{live_v8_html(seed, sat, name, gen_hint)}. The released v7 document
and digest are accepted only when the anchor sat equals
\code{GRANDFATHERED_SAT} (Section~\ref{sec:grandfather}). Version 6 bodies
(\code{0x06}) and every other v7 mint are rejected.
\end{sloppypar}

The fold source must not declare \code{var V8} in the same scope as the
document's \code{V8={...}}. The v6.0 hoisting bug is the reason the document
assigns the object first and the fold only reads it, never re-declares it.

\subsection{The crossing body}

The crossing body is exactly 55 bytes.

\begin{table}[htbp]
\centering
\small
\caption{The \code{FRG8CROSS} state body. The legacy \code{FRG7CROSS}
spelling is accepted only for the grandfathered sat.}
\label{tab:state-body}
\begin{tabular}{@{}rrll@{}}
\toprule
\textbf{Offset} & \textbf{Length} & \textbf{Field} & \textbf{Encoding} \\
\midrule
0  & 9  & magic & \code{"FRG8CROSS"} (legacy \code{"FRG7CROSS"}) \\
9  & 1  & version & \code{0x08} (legacy \code{0x07}) \\
10 & 8  & sat & \code{u64} little-endian \\
18 & 32 & mint txid & internal byte order \\
50 & 4  & mint index & \code{u32} little-endian \\
54 & 1  & dir & $0$ left, $1$ straight, $2$ right \\
\bottomrule
\end{tabular}
\end{table}

$(\mathrm{sat},\mathit{mint\_txid},\mathit{mint\_index})$ identify the frog
and are immutable for its life. \code{dir} is per crossing: it names the
direction the player committed to and therefore the leaf that was spent
(Section~\ref{sec:custody}). A body with \code{dir > 2} is malformed and
rejected.

The mint inscription id
\[
\text{\code{{reversed_mint_txid}i{mint_index}}},
\]
is the ordinal parent of every crossing inscription. The body stores the
same identity in binary, with the txid in internal order. A crossing body
therefore binds four things at once: the sat, the mint transaction, the mint
output and the chosen direction. A witness copied from another frog cannot
be replayed, because the body would not match.

\subsection{The fold algorithm}

The fold is the canonical state reconstruction. It is stateless with respect
to heights: it asks the viewer's ord node for the current tip, walks the
crossing children of the mint, verifies each body-valid child against its raw
transaction, and replays the children in confirmation order. The children
endpoint is followed as \code{{"ids":[...],"more":bool,"page":n}} until
\code{more} is false. A reorganization is handled by asking again; no local
history is authoritative.

For a candidate child $c$ with decoded direction $\mathit{dir}$, the fold
uses three predicates. The first is the body check:
\begin{lstlisting}[language={},caption={Body, raw-transaction and proof predicates.},label={lst:fold-predicates}]
bodyOk(c) = content(c) is exactly 55 bytes,
            starts FRG8CROSS || 0x08 || sat || mint_txid || mint_index,
            has dir <= 2, and the content type is octet-stream

raw(c)    = GET /r/tx/{c.txid}

proofOk(c)= raw(c) is version 2, locktime 0, all input sequences 0xffffffff,
            has exactly one envelope in input 0 and no pointer or extra
            envelope, output 0 has value > 0, the input-0 tapscript is the
            canonical 161-byte leaf for body(c).dir, and its 97-byte control
            block has NUMS as internal key and the canonical sibling path.
            A confirmed Bitcoin transaction is already consensus-valid; the
            fold only has to check that the spent leaf is the canonical one.
\end{lstlisting}
A body-valid child whose proof is missing or false makes the fold
\emph{unverified}; it is never silently skipped. A body-invalid child (wrong
magic, wrong sat, wrong mint, malformed direction) is ignored. Incomplete
pagination, with \code{more} still true after the bounded cursor, is also
\code{unverified}, never a lower-bound count.

\begin{algorithm}[htbp]
\caption{The canonical proof-checking fold}\label{alg:fold}
\begin{algorithmic}[1]
\Require $\mathrm{sat}$, own inscription id, and the embedded parameters
\State $\mathit{id}\gets$ own inscription id from \code{location.pathname}
\State $\mathit{mint}\gets$ parent of $\mathit{id}$ (the mint document is its own mint)
\State $\mathit{tip}\gets \textsc{Get}$\code{(/r/blockheight)}
\State $\mathit{last}\gets \textsc{Get}$\code{(/r/inscription/{id})}
\State $C\gets [\,]$
\For{each page of \textsc{Get}\code{(/r/children/{id})} until \code{more} is false}
  \For{each $c$ in the page's \code{ids}}
    \State $d\gets \textsc{Get}$\code{(/r/inscription/{c})}
    \If{$d.\mathit{sat}=\mathrm{sat}$ and
        $d.\mathit{satpoint}=\mathit{last}.\mathit{satpoint}$ and
        $d.\mathit{height}$ is set}
      \If{$\textsc{BodyOk}(c)$}
        \If{not $\textsc{ProofOk}(c)$}
          \State \Return \textsc{unverified}(\text{``proof''})
        \EndIf
        \State append $(c,\mathit{dir})$ to $C$
      \EndIf
    \EndIf
  \EndFor
\EndFor
\If{\textsc{More} is still true after the bounded cursor}
  \State \Return \textsc{unverified}(\text{``pagination''})
\EndIf
\State sort $C$ by $(\mathit{height},\mathit{sequence\_number})$
\State $\mathit{state}\gets \textsc{idle}$;\quad
       $\mathit{dist}\gets 0$;\quad
       $\mathit{lastT}\gets \text{null}$;\quad $\mathit{txids}\gets \varnothing$
\For{each $(c,\mathit{dir})$ in $C$}
  \If{$c.\mathit{id}\in\mathit{txids}$}
    \State \Return \textsc{unverified}(\text{``duplicate crossing''})
  \EndIf
  \If{$\mathit{lastT}\ne\text{null}$ and
      $c.\mathit{height}-\mathit{lastT}<\kappa$}
    \State \Return \textsc{unverified}(\text{``crossing before finality''})
  \EndIf
  \State $T\gets c.\mathit{height}$
  \If{$T+\kappa>\mathit{tip}$}
    \State \Return \textsc{crossing}$(\mathit{dist},\mathit{dir})$
  \EndIf
  \If{$\mathit{dir}=w(T+1)$}
    \State \Return \textsc{dead}$(\mathit{dist})$
  \EndIf
  \State $\mathit{dist}\gets \mathit{dist}+1$;\quad
         $\mathit{lastT}\gets T$;\quad
         $\mathit{txids}\gets\mathit{txids}\cup\{c.\mathit{id}\}$
  \If{$\mathit{dist}\ge r(T)$}
    \State \Return \textsc{home}$(\mathit{dist})$
  \EndIf
\EndFor
\State \Return \textsc{idle}$(\mathit{dist},\,
       r(\mathit{tip}+1))$
\end{algorithmic}
\end{algorithm}

Three details are consensus-critical. First, a child is admitted only when
its sat and satpoint equal the mint's, its body is the canonical 55-byte body
for this frog and mint, and its raw transaction proves the canonical state
leaf was spent. A parent tag alone is never trusted. Second, the outcome is
recomputed from the block hash and the direction in the body: the fold never
asks the HOPPER indexer what happened; it verifies the spend from ord's raw
transaction endpoint. Third, the finality check rejects a second crossing
before $\kappa$ blocks even when both transactions are already confirmed at
the current tip. A transfer remains invisible to \code{dist}, but it cannot
create a false \code{HOME}: the next crossing lacks a canonical state leaf
and fails the proof.

\subsection{The one grandfathered v7 Hopper}\label{sec:grandfather}

Exactly one released Hopper is carried across the v8 cut. Its sat is the
global \code{GRANDFATHERED_SAT} value:
\begin{itemize}[leftmargin=1.4em,itemsep=0.2em]
  \item Rust: \code{hopper_core::params::GRANDFATHERED_SAT};
  \item JavaScript: \code{GRANDFATHERED_SAT} in
        \code{web/js/core/params.js}.
\end{itemize}
The production value is \code{1_965_983_893_977_665}; \code{0} means none.
It is the only value to edit when the sat changes, and it is embedded into
new v8 mint documents as \code{grandfatheredSat} so the fold does not need a
second copy.

For that sat only:
\begin{itemize}[leftmargin=1.4em,itemsep=0.2em]
  \item the v7 mint envelope (\code{0x07}), the v7 document digest and the
        \code{V7} body are accepted;
  \item the v7 state body (\code{FRG7CROSS}/\code{0x07}) and the v7
        161-byte leaf and control path are accepted;
  \item the wallet uses the released leaf and version and can recover or
        migrate the sat on the legacy Taproot derivation
        (Section~\ref{sec:wallet}).
\end{itemize}
Every other sat must use v8. A v7 crossing for a non-grandfathered sat is
rejected by the indexer and is not counted by a v8 fold; the app and indexer
report it as \code{unverified}. The grandfathered Hopper is not a precedent
for a second legacy frog.

The released v7 card's embedded fold is immutable and predates the
proof-checking algorithm. It is grandfathered for custody and play, not as an
authoritative art oracle; a viewer that needs the proof-checked state must
use the v8 app or indexer, or a new v8 card.

\begin{invariant}{No hidden state}
Everything the fold needs is in the mint document, the crossing bodies, the
raw transactions, the block hashes and the ord node's satpoint answers. There
is no HOPPER database in the trust path, and a body is never believed merely
because it is attached to the right parent.
\end{invariant}
